\section{This Water, This Bread}

\subsection{The particularity did not retire at the Ascension}

A tempting compromise would concede everything so far and then quietly re-abstract the result: granted, God acted particularly \emph{then}; but now the Incarnation survives as a truth to be believed, an example to be admired---portable, general, safely conceptual. The Church's sacramental teaching exists to block precisely that retreat. Leo the Great, preaching on the Ascension, states the principle in a single clause: ``that which till then was visible of our Redeemer was changed into a sacramental presence.''\endnote{Leo the Great, \emph{Sermo} 74 (De ascensione Domini II), 2, in the public-domain Nicene and Post-Nicene Fathers translation (Charles Lett Feltoe), checked at the New Advent transcription, 24 July 2026.} The visible body departs; the mode of presence changes; the particularity does not. What replaces the carpenter's hands is not an idea but a set of instituted, perceptible acts---washing, anointing, the giving of bread and the cup---each as concrete, as datable, as located as Nazareth.

The Catechism explains why this is not a lapse back into magic but the continuation of the same divine style. The sacraments are ``perceptible signs (words and actions) accessible to our human nature,'' by which the risen Christ, seated at the Father's right hand, now acts to communicate the grace they signify. And they are possible because the Paschal event is unlike every other event: all others ``happen once, and then they pass away, swallowed up in the past,'' but the Paschal mystery of Christ ``cannot remain only in the past,'' since in it death itself was destroyed and everything Christ did and suffered participates in divine eternity.\endnote{\emph{Catechism of the Catholic Church} 1084--1085, official English text at the Vatican website, checked 24 July 2026. This is a claim of sacramental making-present, not of repetition: the historical event is not reenacted but its saving power made available. The article does not expand here into a treatise on sacramental validity, matter and form, or Eucharistic presence, which have their own governing sources.} Lessing's ditch was dug on the assumption that a past event can reach the present only by report, dwindling as testimony dwindles. The sacramental economy is the claim that this one event reaches the present by \emph{presence}---that the ditch is crossed, not by our leap backward, but by the event's own refusal to recede.

\subsection{Why signs suit us}

Aquinas's defense of sensible signs completes the pattern this article has been tracing since the discussion of abstraction. Divine wisdom, he says, ``provides for each thing according to its mode,'' and it is the human mode---the point was established about knowers, and returns here about worshipers---``to acquire knowledge of the intelligible from the sensible.'' Hence sacraments are sensible things; and not sensible things in general, but ``those things which are determined by Divine institution.''\endnote{Aquinas, \emph{Summa theologiae} III, q.~60, aa.~4--5, English Dominican translation, checked 24 July 2026.} A purely spiritual religion---all interiority, no matter---would be fitted for angels, and its persistent attraction for us is one more instance of the flight into abstraction that this article has been resisting throughout: the preference for a God who makes no landings. The actual economy is humbler and more invasive. Grace comes to a creature of flesh through flesh: through water that is wet, bread that is bread, words said aloud by one particular minister to one particular person with a name.

For the recipient, this particularity is not a limit on assurance but its instrument. An abstract salvation must be verified abstractly---by introspection, by the quality of one's feelings, a notoriously bottomless audit. A sacrament happens or it does not. It happened to \emph{you}, on a date, at a font one can drive to. The definite articles that scandalized Nazareth turn out, at the receiving end, to be the whole comfort: not humanity in general invited to a possible mercy, but this person, washed on this morning, in this parish, by name.

The universal reach that the objector feared was being lost is, once again, achieved rather than abandoned. Precisely because the sacramental acts are small, repeatable, and material, they can be performed anywhere humans are---in cathedrals and prison cells, over any water, in every language---until the particular address has been multiplied toward every particular person. A universal offer that remained abstract would reach no one; the concrete sign, endlessly reiterated, is how the once-for-all event goes out ``even to the farthest part of the earth.''
